\section{Rangkuman}
\label{sec:rangkuman-bab2}

Pada bab ini telah dibahas:

\begin{itemize}
    \item Model mesin von Neumann terdiri atas unit pemroses, memori,
          unit masukan, dan unit keluaran. Ciri utamanya adalah
          prinsip \eng{stored-program}: instruksi dan data disimpan
          di memori yang sama.

    \item Siklus instruksi terdiri atas \eng{fetch}, \eng{decode},
          dan \eng{execute}, dengan \eng{store} sebagai langkah
          tambahan pada 8086. Satu instruksi dapat memerlukan
          beberapa siklus mesin.

    \item ISA adalah kontrak instruksi antara perangkat lunak dan
          perangkat keras. Kompatibilitas ISA membuat program 8086
          masih dapat dijalankan pada prosesor x86 modern.

    \item \eng{Opcode} menyatakan operasi, \eng{operand} menyatakan
          data, dan \eng{mnemonic} adalah nama pendek opcode dalam
          bahasa assembly.

    \item Register adalah penyimpan kecil di dalam prosesor dengan
          akses jauh lebih cepat daripada memori. Jumlahnya terbatas,
          sehingga harus dikelola secara sadar.

    \item Penyandian data di level mesin meliputi biner,
          heksadesimal, \eng{two's complement}, BCD, dan ASCII.
          Representasi yang sama dapat memiliki tafsiran berbeda
          bergantung pada apakah bilangan diperlakukan bertanda atau
          tak bertanda.

    \item \eng{Carry} (\flag{CF}) menandakan kelebihan pada operasi
          tak bertanda, sedangkan \eng{overflow} (\flag{OF})
          menandakan kelebihan pada operasi bertanda. Keduanya harus
          dibedakan.

    \item Peta konsep buku menunjukkan empat klaster kompetensi yang
          saling bergantung, dan urutan bab mengikuti
          ketergantungan tersebut.
\end{itemize}

\begin{checklistbox}
Tandai setiap butir yang sudah Anda kuasai.

\begin{itemize}
    \item[\checkbox] Saya dapat menjelaskan model von Neumann dan
          prinsip \eng{stored-program}.
    \item[\checkbox] Saya dapat menjelaskan tahapan siklus instruksi.
    \item[\checkbox] Saya dapat mendefinisikan ISA, \eng{opcode},
          \eng{operand}, \eng{mnemonic}, dan register dengan kalimat
          saya sendiri.
    \item[\checkbox] Saya dapat mengonversi bilangan antara biner,
          heksadesimal, dan desimal tanpa alat bantu.
    \item[\checkbox] Saya dapat menyatakan bilangan negatif dalam
          \eng{two's complement} dan memverifikasinya.
    \item[\checkbox] Saya dapat membedakan \eng{carry} dan
          \eng{overflow} beserta flag yang menandainya.
    \item[\checkbox] Saya dapat menjelaskan penyandian BCD dan ASCII
          serta perbedaan nilai dengan sandian karakter.
    \item[\checkbox] Saya dapat membaca peta konsep buku dan
          menjelaskan urutan ketergantungan antarbab.
\end{itemize}

\vspace{0.4em}
Bila masih ada butir yang belum bertanda, ulangi bagian yang
bersesuaian sebelum melanjutkan ke Bab~\ref{chap:pengenalan}.
\end{checklistbox}

\begin{infobox}
\textbf{Rujukan bab.} Materi pada bab ini diperkuat oleh
\cite{tanenbaum2013}, \cite{stallings2016}, dan \cite{brey2009}
untuk organisasi komputer dan representasi data. Untuk gambaran
teknis 8086 yang lebih rinci, lihat \cite{intel1979}. Pembahasan
\eng{two's complement} dan operasi bertanda pada bahasa assembly
dapat dilihat pada \cite{hyde2010}.
\end{infobox}
